\chapter{Spiele erstellen}
\label{ch:games/games}

Die Stärke von Luanti liegt darin, dass sich Spiele einfach entwickeln
lassen, ohne eigene Voxel-Grafik, Voxel-Algorithmen oder ausgefeilten
Netzwerkcode schreiben zu müssen.

\section{Was ist ein Spiel?}\label{was-ist-ein-spiel}

Ein Spiel ist eine Sammlung von Mods, die zusammen ein stimmiges Ganzes
ergeben. Ein gutes Spiel hat ein durchgängiges Thema und eine Richtung:
Es kann zum Beispiel ein klassisches Bau- und Abbauspiel mit harten
Überlebenselementen sein oder eine Weltraumsimulation mit
Steampunk-Automatisierung.

Spieldesign ist ein komplexes Thema und ein eigenes Fachgebiet. Dieses
Buch kann es nur kurz streifen.

\section{Spielverzeichnis}\label{spielverzeichnis}

Aufbau und Ort eines Spiels kommen dir nach der Arbeit mit Mods recht
vertraut vor. Spiele liegen in einem Spielverzeichnis, etwa
\texttt{minetest/\allowbreak{}games/\allowbreak{}foo\_\allowbreak{}game}.

\begin{Shaded}
\begin{Highlighting}[]
\NormalTok{foo\_game}
\NormalTok{├── game.conf}
\NormalTok{├── menu}
\NormalTok{│   ├── header.png}
\NormalTok{│   ├── background.png}
\NormalTok{│   └── icon.png}
\NormalTok{├── minetest.conf}
\NormalTok{├── mods}
\NormalTok{│   └── ... Mods}
\NormalTok{├── README.txt}
\NormalTok{└── settingtypes.txt}
\end{Highlighting}
\end{Shaded}

Nötig ist nur der Ordner \texttt{mods};
\texttt{game.\allowbreak{}conf}\index{game.conf@\texttt{game.\allowbreak{}conf}}
und \texttt{menu/\allowbreak{}icon.\allowbreak{}png} werden aber
empfohlen. Eine einfache \texttt{game.\allowbreak{}conf} sieht so aus:

\begin{Shaded}
\begin{Highlighting}[]
\NormalTok{title = Foo Game}
\NormalTok{description = Ein Spiel über Foo.}
\end{Highlighting}
\end{Shaded}

\section{Kompatibilität zwischen
Spielen}\label{kompatibilituxe4t-zwischen-spielen}

\subsection{API-Kompatibilität}\label{api-kompatibilituxe4t}

Es lohnt sich, so viel API-Kompatibilität mit Minetest Game zu wahren,
wie es sich einrichten lässt. Dann lassen sich Mods viel leichter auf
ein anderes Spiel übertragen.

Am besten wahrst du die Kompatibilität mit einem anderen Spiel, indem du
die API-Kompatibilität mit gleichnamigen Mods wahrst. Trägt ein Mod
denselben Namen wie ein anderer -- auch wie ein Mod von Dritten --,
sollte er also eine kompatible API haben. Enthält ein Spiel zum Beispiel
einen Mod namens \texttt{doors}, sollte er dieselbe API haben wie
\texttt{doors} in Minetest Game.

Die API-Kompatibilität eines Mods setzt sich zusammen aus:

\begin{itemize}
\tightlist
\item
  der Lua-API-Tabelle -- alle dokumentierten oder beworbenen Funktionen
  in der gleichnamigen globalen Tabelle, zum Beispiel
  \texttt{mobs.\allowbreak{}register\_\allowbreak{}mob};
\item
  den registrierten Nodes und Items -- also dass diese vorhanden sind.
\end{itemize}

Kleine Brüche sind nicht so schlimm, etwa wenn eine beiläufige
Hilfsfunktion fehlt, die eigentlich nur intern verwendet wurde. Größere
Brüche bei Kernfunktionen sind dagegen sehr schlecht.

Mit einem riesigen Alleskönner-Mod wie \emph{default} in Minetest Game
ist API-Kompatibilität kaum zu wahren. In diesem Fall sollte das Spiel
keinen Mod namens \texttt{default} enthalten.

API-Kompatibilität betrifft auch andere Mods und Spiele von Dritten.
Achte also darauf, dass neue Mods einen eindeutigen Namen haben. Ob ein
Mod-Name schon vergeben ist, findest du mit einer Suche auf
\url{https://content.luanti.org/} heraus.

\subsection{Gruppen und Aliase}\label{gruppen-und-aliase}

Gruppen\index{gruppe@Gruppe} und Aliase\index{alias@Alias} helfen beide,
die Kompatibilität zwischen Spielen zu wahren, denn mit ihnen dürfen
sich Itemnamen zwischen Spielen unterscheiden. Häufige Nodes wie Stein
und Holz sollten Gruppen haben, die das Material angeben. Es ist
außerdem sinnvoll, Aliase von \texttt{default}-Nodes auf ihre direkten
Entsprechungen anzulegen.

\section{Jetzt bist du dran}\label{jetzt-bist-du-dran}

\begin{itemize}
\tightlist
\item
  Erstelle ein einfaches Spiel, in dem der Spieler Punkte bekommt, wenn
  er besondere Blöcke abbaut.
\end{itemize}
